\section{Fox colorings and finite braid actions}
\label{sec:fox}

A coloring gives a short witness that a knot is nontrivial, and counting
colorings can give the exact number of crossing changes needed to untie a
whole family. We work throughout with finite diagrams and explicit move
sequences. The classical names for the braid closures below specify their
usual diagrams. Background sources for this section are
J.~H.~Przytycki, \emph{3-coloring and other elementary invariants of knots}
(\href{https://arxiv.org/abs/math/0608172}{arXiv:math/0608172}),
L.~H.~Kauffman and P.~Lopes, \emph{Determinants of rational knots}
(\href{https://dmtcs.episciences.org/articles/454}{DMTCS}), and
S.~Asami and K.~Kuga, \emph{Colorings of torus knots and their
twist-spins by Alexander quandles over finite fields}
(\href{https://www.math.s.chiba-u.ac.jp/~xasami/library/coloring.pdf}{author's paper}).

\subsection{A local rule that survives every move}

Fix a modulus $n\ge2$ and put $R=\mathbb Z/n\mathbb Z$, implemented by
residues $0,\ldots,n-1$. Cut the diagram at every crossing into edge
segments, giving each crossing-free component a single segment. At a
crossing require equality of the two overpassing segment colors, say $y$,
and require the two underpassing segment colors $x,z$ to satisfy
$x+z=2y$. Colors therefore propagate by the reflection
$S_y(x)=2y-x$. Write $C_n(D)$ for the set of assignments satisfying all
these equations and $F_n(D)$ for its cardinality. Equivalently one may use
one variable per arc between successive undercrossings, including a single
variable for each component with no undercrossing.

\begin{theorem}[Coloring invariance and tricolorability]
Each Reidemeister move induces a bijection of Fox colorings, fixed on the
exterior of the move disc. These bijections preserve constant colorings.
A knot is tricolorable exactly when $F_3(D)>3$, and every such knot coloring
which is nonconstant uses all three colors.
\end{theorem}
\begin{proof}
The identities needed at the three moves are
\[
 S_x(x)=x,\qquad S_y(S_y(x))=x,\qquad
 S_z(S_y(x))=S_{S_z(y)}(S_z(x)).
\]
The last identity has value $2z-2y+x$ on both sides. At a kink the strand
color uniquely extends to the little loop by idempotence and invertibility.
In type II the overstrand has color $y$ and the intervening understrand
must have color $S_y(x)$, so its exit has color $x$. In type III name the
bottom, middle and top colors on entry $x,y,z$. The two bottom exits are
the two sides of the last identity. The middle exit is $S_z(y)$ on both
sides and the top exit is $z$. Each interior segment has been determined
by the entry data. These descriptions apply from either end of an
understrand since $S_y$ is its own inverse. Relabeling and reversing the
reading of the local picture give the other presentations of the moves.
Exterior assignments therefore extend to the two sides in pairs.
For $n=3$ the equation reads $x+y+z=0$. Two equal entries force all
three equal, and the three distinct residues satisfy the equation. There
are three constant assignments. A nonconstant assignment on a single
component must change color at some underpass, since otherwise traversal
would give a single color everywhere. That crossing uses all three.
\end{proof}

The coloring equations are homogeneous and linear. Hence coordinatewise
addition, negation, and scalar multiplication preserve solutions. On a
nonempty diagram with marked segment $s$, translation gives a bijection
\[
 C_n(D)\longleftrightarrow R\times C_n^0(D,s),\qquad
 f\longmapsto(f(s),f-f(s)),
\]
where $C_n^0$ means $f(s)=0$ and the second term subtracts the constant
coloring. Disjoint unions multiply the counts, so the $k$-component
unlink has $n^k$ colorings. Modulo $2$ the understrand keeps its original
color, hence any $k$-component diagram has precisely $2^k$ colorings.
For coprime $a,b\ge2$, $F_{ab}(D)=F_a(D)F_b(D)$: reduction of every
entry is a bijection by the Chinese remainder construction and respects
each equation in both directions. Concretely, given residues $u,v$ use
$u+a((v-u)a^{-1}\bmod b)$ with $a^{-1}$ obtained by bounded search.

For a prime $\ell$ the count is always a power of $\ell$. To see this
constructively, let $W_j$ be the set of colorings vanishing on the first
$j$ segments, $0\le j\le e$, where $e$ is the total number of segments.
Look for an element of $W_j$ whose next entry is nonzero. If the bounded
search fails, $W_j=W_{j+1}$. If it succeeds, rescale that element to
obtain $v$ with next entry $1$. Then $(a,w)\mapsto av+w$ is a bijection
from $(\mathbb Z/\ell)\times W_{j+1}$ onto $W_j$; its inverse reads off
the next entry $a$ and subtracts $av$. Since $|W_e|=1$, scanning backward
gives $F_\ell(D)=\ell^k$, where $k$ counts the successful searches.

\subsection{Braids, fixed points and torus knots}

Read a braid from top to bottom with strands numbered from the left.
A positive generator takes input colors $(x,y)$ with the left strand
overpassing to $P(x,y)=(2x-y,x)$, and its inverse takes $(u,v)$ to
$(v,2v-u)$. Write $P_i$ for this update on neighboring coordinates $i,i+1$.
Coloring a braid is deterministic given its top row. Closure joins equal
positions, so its colorings are in bijection with fixed rows of the word.
In particular a closed $b$-strand braid has at most $n^b$ colorings.
This representation respects inverse cancellation and commutation at
disjoint positions. On three coordinates both $P_1P_2P_1$ and $P_2P_1P_2$
give $(2x-2y+z,2x-y,x)$, proving the braid relation. A conjugation of
the word carries fixed rows bijectively by the conjugating permutation.
For a stabilization, write the last output of the original word as $u$
and the new input as $v$. The positive final crossing gives
$(2u-v,u)$, the negative one gives $(v,2v-u)$. In each case the closure
equations force $u=v$ and then require that the original last input was $u$.
This proves invariance under both braid stabilizations directly.

\begin{proposition}[All two-strand powers]
For every signed integer $m$, the closure of $\sigma_1^m$ has
$n\gcd(n,|m|)$ colorings, including $n^2$ at $m=0$. Its component count
is one for odd $m$ and two for even $m$.
\end{proposition}
\begin{proof}
Put $d=x-y$. Then $P(x,y)=(x+d,y+d)$ preserves $d$. Induction, applying
$P^{-1}$ for negative powers, gives $(x+md,y+md)$. The row closes exactly
when $md=0$ in $R$. For $g=\gcd(n,|m|)$ this says that $d$ is a multiple
of $n/g$: after cancellation modulo $n/g$ the number $|m|/g$ is a unit
when $n/g>1$, by Euclid's algorithm. The $g$ representatives
$j(n/g)$, $0\le j<g$, and the $n$ choices for $x$ specify all solutions
uniquely. The endpoint permutation is the transposition to the $m$th
power, proving the component count.
\end{proof}

For the trefoil ($m=3$) there are nine tricolorings, versus three for the
unknot. For the Hopf link ($m=2$) there are three, versus nine for the
two-component unlink. More generally the knots with odd $m=1,3,5,\ldots$
are pairwise inequivalent. For positive odd $a<b$, take $n=b$.
The larger diagram has $b^2$ colorings and the smaller has at most $ab<b^2$.
This also furnishes a primitive recursive sequence of arbitrarily many
distinct knot types in the finite calculus.

\begin{theorem}[All coprime positive torus parameters]
Let $p,q\ge1$ be coprime and let $T(p,q)$ be the standard closure of
$(\sigma_1\cdots\sigma_{p-1})^q$. Its Fox count for every modulus $n\ge2$ is
\[
 F_n(T(p,q))=
 \begin{cases}
 n & p,q\text{ both odd},\\
 n\gcd(n,p)&q\text{ even},\\
 n\gcd(n,q)&p\text{ even}.
 \end{cases}
\]
More generally, for arbitrary positive $p,q$ the standard torus link has
$\gcd(p,q)$ components.
\end{theorem}
\begin{proof}
The permutation of a single block is a cyclic shift by one. Closure
therefore follows addition by $q$ on $\mathbb Z/p\mathbb Z$ (the opposite
shift has the same orbits). Its orbits are the residue classes modulo
$\gcd(p,q)$, by cancellation after dividing by that gcd.

Use cyclic indices $i=0,\ldots,p-1$, and set $d_i=x_{i+1}-x_i$. The
first strand passes over all the others. A block $B$ sends $x_i$ to
$(Bx)_i=2x_0-x_{i+1}$, including $(Bx)_{p-1}=x_0$. Hence it sends
$d_i$ to $-d_{i+1}$. A row fixed by $B^q$ satisfies
$d_i=(-1)^q d_{i+q}$ and $\sum_i d_i=0$. For coprime $p,q$ the shift by
$q$ visits all positions.

If $q$ is even all differences equal a single $d$, with $pd=0$.
Conversely the row $x_i=x_0+id$ with $pd=0$ changes under one block to
the row starting at $x_0-d$ and having difference $-d$. A second block
restores it. This gives exactly $n\gcd(n,p)$ fixed rows.
If $p,q$ are odd, iterating the signed shift $p$ times gives $d_i=-d_i$.
All signs are then immaterial, so all differences equal $d$ with
$2d=pd=0$. As $p$ is odd, $d=0$, and the constant rows give exactly $n$
solutions. Finally, if $p$ is even, $q$ is odd and each shift changes
parity of the index. Thus $d_i=(-1)^i d$ for some $d$, and the row
alternates between $x_0$ and $x_0+d$. One block subtracts $d$ from
each entry, preserving all its differences. Closure is exactly $qd=0$.
The two-strand counting argument gives $n\gcd(n,q)$.
\end{proof}

The standard figure-eight diagram is the closure of
$(\sigma_1\sigma_2^{-1})^2$. A single block sends $(x,y,z)$ to
$(2x-y,z,2z-x)$. Put $y=x+d$ and $z=x+e$. On applying two blocks,
equality of the three output positions to their inputs gives, respectively,
$2d+e=0$, $-d+2e=0$, and $d+3e=0$. These equations are equivalent to
$e=-2d,\ 5d=0$, so the figure-eight has exactly $n\gcd(n,5)$ colorings.
Modulus $3$ distinguishes it from the trefoil, and modulus $5$ distinguishes
it from the unknot.

\subsection{Conway rational tangles by two boundary differences}

Number a rectangular two-string tangle's boundary colors clockwise,
starting at the northwest corner, by $a,b,c,d$. Write
$h=b-a$ and $v=d-a$. Begin with the horizontal zero tangle whose
northern pair and southern pair are each joined. Generate rational
tangles by right twists, inverse right twists and clockwise quarter-turns.
For the right twist use the crossing in which the old southeast strand
goes over the old northeast one to become the new northeast strand.
This fixes the fraction convention, including its sign.

\begin{theorem}[Color count of both rational closures]
Start the integer pair at $(P,Q)=(0,1)$. A right twist replaces it by
$(P+Q,Q)$, an inverse twist by $(P-Q,Q)$ and a quarter-turn by
$(-Q,P)$. After any word of operations, joining the northern and southern
pairs (numerator closure) gives $n\gcd(n,|P|)$ Fox colorings.
Joining the western and eastern pairs (denominator closure) gives
$n\gcd(n,|Q|)$. The numerator closure is a knot exactly when $P$ is odd,
and a numerator knot is tricolorable exactly when $3\mid P$.
\end{theorem}
\begin{proof}
We maintain a bijection between all colorings of the tangle and pairs
$(a,u)\in R^2$, with boundary values
\[
 (a,b,c,d)=(a,\ a+Pu,\ a+(P+Q)u,\ a+Qu).
\]
It holds at the start since the northern color is $a$ and the southern
color can be any $a+u$. A twist replaces the eastern pair $(b,c)$ by
$(c,2c-b)$ and fixes the western pair; the old assignment and the
extended one determine each other. The new boundary formula has
$P'=P+Q$, $Q'=Q$. The inverse twist reverses this substitution.
Rotation changes $(a,b,c,d)$ to $(d,a,b,c)$. The new northwest parameter
is $a'=a+Qu$ and $u$ is unchanged. The new differences are $-Qu$ and
$Pu$, giving the asserted pair and a reversible parameter change.
This proves the bijection by induction on the construction word.
In particular the integer pairs remain coprime, since each update
preserves their gcd. Their slopes follow the familiar operations
$P/Q\mapsto P/Q\pm1$ and $P/Q\mapsto -Q/P$, with the integer pair
also covering a zero denominator.

The two numerator joins impose $a=b$ and $c=d$, both equivalent to
$Pu=0$. The denominator joins impose $a=d$ and $b=c$, both equivalent
to $Qu=0$. Counting $u$ as in the two-strand proof yields both formulas.
At $n=2$ the general component count $2^k$ gives one component for
odd $P$ and two for even $P$. At $n=3$ compare the numerator count with
three constant colorings. At each construction step $|P|+|Q|$ at most
doubles, so $2^L$ bounds it after a word of length $L$.
\end{proof}

For example, from $(0,1)$ apply two right twists, turn clockwise, and
apply three more right twists. This gives $(2,1)$, then $(-1,2)$, then
$(5,2)$, so the numerator is a knot with $n\gcd(n,5)$ colorings.
The formula uses only a designated construction word. Classification
of rational tangles by their fractions would assert a further converse.

\subsection{A finite invariant beyond every Fox modulus}

On an oriented diagram one may replace $S_y$ by the affine permutation
$S_y(x)=t x+(1-t)y$ for any unit $t$ modulo $n$. At each oriented overpass
use its normal to order the two underpassing sides; the two under-colors
in this order are $x,S_y(x)$. Choose the normal convention for which a
positive downward braid generator gives $(S_x(y),x)$. When traversing
against this normal apply the inverse $S_y^{-1}(x)=t^{-1}x+(1-t^{-1})y$.
Again $S_y(y)=y$. Every affine map $A(x)=ax+b$ satisfies
$A S_y(x)=S_{A(y)} A(x)$ by expansion. In particular this holds for
$A=S_z$ and $S_z^{-1}$ and with $S_y$ replaced by its inverse (when $A$
is invertible, conjugate the original identity and take inverses).
These identities give all oriented type III relations: moving across
the top strand transports both the middle color and its action by
$A$, in the direction prescribed by its normal. Type II uses mutually
inverse actions and type I uses the fixed point $y$. The extension and
restriction argument above therefore proves oriented move invariance for
all these counts.

For $T(3,5)$ every Fox modulus gives the unknot's count. Nevertheless, a
single explicit affine coloring distinguishes this knot. Use $n=31$,
$t=7$ (inverse $9$). Successive rows after each full positive
$\sigma_1\sigma_2$ block are
\[
\begin{aligned}
(0,1,6)&\longmapsto(7,11,0)\longmapsto(4,20,7)\\
 &\longmapsto(23,25,4)\longmapsto(6,14,23)\longmapsto(0,1,6).
\end{aligned}
\]
The row closes and is nonconstant. Its translates and nonzero scalar
multiples are also fixed. In fact $a+b(0,1,6)$ for $a,b\in R$ gives
$31^2$ distinct fixed rows by the first two coordinates. This count is exact.
The first coordinates of the fifth block on the three unit rows, evaluated
by the same recurrence, are $13,29,21$. A fixed row therefore satisfies
$12x+29y+21z=0$, which forces $z=-5x+6y$ since $21^{-1}=3$ modulo $31$.
The crossing-free circle has only $31$. An unoriented equivalence to the unknot would
transport the given orientation along the certificate to one of the
two orientations of that circle, both having the same $31$ colorings.

\subsection{Connected sum, bridges and crossing changes}

A convenient finite presentation of a $(1,1)$-tangle is a slice word in a
rectangle, with one strand end at the top and one at the bottom. An event
is a crossing of adjacent strands, the insertion of two adjacent ends of
a local extremum, or the deletion of two adjacent ends meeting at an
extremum. Each cup or cap joins equal colors. For a polygonal diagram
in general position this word is read by sorting all vertex and crossing
heights; an ordinary bend preserves color. Cut a small arc on the outside
of a planar knot diagram and run its two ends along the exterior to the
top and bottom of a containing rectangle. Joining the two ends outside
recovers its closure. Equivalently use this marked tangle as input from
the start. The connected sum of two such marked knot closures is the
closure of their vertical concatenation.

\begin{lemma}[Endpoint equality]
Every Fox coloring of a $(1,1)$-tangle assigns the same color to its two
ends, over every $R$ above.
\end{lemma}
\begin{proof}
Read the colors at a regular horizontal slice from left to right, and form
$x_1-x_2+x_3-\cdots$. A crossing preserves the alternating sum since
$(2x-y)-x=x-y$ and the inverse preserves it too. A cup inserts adjacent
entries $u,u$, whose contributions cancel; entries further right keep
their original signs. A cap has the inverse effect. The sums on the top
and bottom slices each consist of a single term, so the end colors agree.
\end{proof}

\begin{theorem}[Multiplication under connected sum]
For the marked knot diagrams just described,
\[
 n F_n(D\mathbin{\#}E)=F_n(D)F_n(E).
\]
In particular the reduced count $F_n(D)/n$ is multiplicative.
\end{theorem}
\begin{proof}
Endpoint equality identifies colorings of each open tangle with those of
its closure: the exterior joining arc has its uniquely prescribed color.
A coloring of the concatenation is uniquely a pair of such colorings
agreeing on the joined end. For every fixed end color the number of
choices in each factor is its reduced count by translation. Summing the
product over the $n$ end colors gives the formula. The endpoint argument
is essential here, since cutting open a general system of equations can
admit new solutions.
\end{proof}

\begin{proposition}[Local cost of a crossing change]
For diagrams $D,D'$ differing by one crossing change,
$F_n(D)\le nF_n(D')$ and $F_n(D')\le nF_n(D)$.
\end{proposition}
\begin{proof}
Subdivide the four incident edges at the boundary of the crossing disc
so the two coloring problems use the same exterior variables. On either
side, requiring all four colors there to coincide defines the same set
$H$ of assignments. In $C_n(D)$ take the difference between one
overpassing and one underpassing color. This is a homomorphism to $R$,
with kernel $H$: if the two agree, the other underpassing color is forced
to agree as well. Every nonempty fiber is a translate of $H$ by subtraction
of its first assignment in lexicographic order. Thus
$|H|\le F_n(D)\le n|H|$. The switched crossing has the same bounds.
For prime $n$, the image of either difference map is zero or all of $R$,
as multiplication by a nonzero scalar is bijective. Thus the counts
across a crossing change are equal or differ by exactly a factor $n$.
\end{proof}

\begin{corollary}[A sharp unknotting bound]
The standard connected sum of $r$ trefoils, each of either handedness, has
$3^{r+1}$ tricolorings and unknotting number exactly $r$ in the diagram
calculus. Every equivalent closed braid has at least $r+1$ strands.
Every equivalent closed slice diagram has at least $r+1$ creating extrema
and at least $r+1$ annihilating extrema.
\end{corollary}
\begin{proof}
For $r=0$ take the unknot. The two-strand formula gives nine colorings
for either trefoil, and the connected sum formula multiplies by three at
each insertion. After $h$ crossing changes, interspersed with arbitrary
Reidemeister moves, a diagram equivalent to the unknot could have started
with at most $3^{h+1}$ colorings. Hence $h\ge r$. Conversely switch the
middle crossing in each standard $\sigma_1^{\pm3}$ block. The adjacent
opposite pair cancels by type II, leaving the closure of a one-crossing
two-strand braid, which is a curl removable by type I. These moves can
be made in each summand disc, choosing the cut on the exterior return arc
away from the curl. The resulting concatenation of straight marked
unknots is the unknot.

The closed braid bound follows from $F_3\le3^b$. For a sweep from zero
strands at the bottom back to zero at the top, choosing a color at
each creating extremum determines both newly incident entries. Every
crossing acts deterministically, and a deletion merely checks equality.
Thus the number of accepted colorings is at most $3^u$ for $u$ creating
extrema. Creation and deletion each change the number of strands by two,
so their totals are equal.
\end{proof}

\paragraph{Primitive recursive and Lean details.}
For $e$ segments, enumerate row codes $0,\ldots,n^e-1$ in base $n$ and
test the crossing equations by a bounded conjunction. A braid on $b$
strands acts on a row code below $n^b$ by replacing two digits and
reducing them modulo $n$, iterating once per letter. Comparisons of
functions on these codes are finite Boolean tables. For each local
bijection the finite counting lemma can be proved by summing the
indicator of its graph, once by rows and once by columns, using induction
on partial rectangular sums. To verify the row after $j$ braid blocks,
use the single bounded Boolean conjunction of the asserted coordinate
equalities over $i<p$. The orbit proof needs just $p$ steps, and the
coprime shift has inverse given by a bounded lookup modulo $p$ (the
one-element case is immediate). For a slice word of length $L$ and initial
width $b$, no slice exceeds $b+2L$. The count bound can equivalently be
checked by mapping an assignment to its list of creating-extremum
colors and reconstructing the slice row at each prefix. For a change
sequence the induction assertion is the numerical inequality at the
current prefix, with the power of $n$ indexed by its number of changes.
All the running equalities, inequalities, digit tests and bounded
conjunctions above are quantifier-free after characteristic functions
are named. No minimization over all diagrams is used by the sharp
unknotting claim: it gives a witness of cost $r$ and an inequality for
every supplied certificate.

In Lean one can use natural-number row codes throughout. Define
\texttt{digit}, \texttt{putDigit}, \texttt{step}, \texttt{run},
\texttt{validColor}, and \texttt{count} in that order. Prove first that
putting a digit updates that coordinate and preserves the others; next
prove the three local identities by modular arithmetic, the closure
bijection by induction on the word, and the torus formula by the cyclic
difference lemma. Encode a local move with its exterior identification;
extend colors by the formulas above and prove the inverse identities
pointwise. For connected sums build a slice-word type carrying arities,
then prove preservation of the alternating sum for each constructor.
The crossing-change proof uses a bounded least representative in each
nonempty difference fiber. These choices give executable witnesses for
all the bijections used in this section.
